% CLAUDE REPORT::B3-FRONTIER-0001::2026-11-15T10:52-08:00
\documentclass[11pt]{article}
\usepackage{amsmath,amssymb,amsthm}
\usepackage[margin=1in]{geometry}
\usepackage{hyperref}
\newtheorem{theorem}{Theorem}
\newtheorem{proposition}[theorem]{Proposition}
\newtheorem{corollary}[theorem]{Corollary}
\theoremstyle{definition}
\newtheorem{definition}[theorem]{Definition}
\theoremstyle{remark}
\newtheorem{remark}[theorem]{Remark}
\title{CLAUDE B3-FRONTIER-0001:\\ Where the conditional bridge $\mathrm{SAT}\in\mathrm{P}\Rightarrow \overline W_R(N)\le\mathrm{poly}(N)$ can and cannot carry content}
\author{Claude (scheduled worst-case derivation-cost routine)}
\date{2026-11-15, 10:52 America/Los\_Angeles}
\begin{document}
\maketitle

\noindent\textbf{STATUS SUMMARY.}
Propositions~\ref{closure}--\ref{world}: PROVED/EXACT (elementary).
Theorem~\ref{pvbridge}: KNOWN THEOREM (Cook 1975 propositional translation; Cook--Kraj\'{\i}\v{c}ek 2007), restated in project notation.
Corollary~\ref{split}: PROVED/EXACT given Theorem~\ref{pvbridge}.
Residual gap $\mathrm{RG}(\mathrm{EF})$, any superpolynomial lower bound on $\overline W_{\mathrm{EF}}$, and P versus NP: \textbf{OPEN}.
This note proves no lower bound and no unconditional bridge. P versus NP remains OPEN.

\section{Setup (as in the 2026-11-02 worst-case directive)}
Fix a sound, complete Cook--Reckhow refutation system $R$ for CNF (proofs checkable in polynomial time). For UNSAT $F$,
\[
W_R(F)=\min_{\pi:F\vdash_R\bot}\operatorname{cost}_R(\pi),\qquad
\overline W_R(N)=\max_{\substack{F\in\mathrm{UNSAT}\\ |\langle F\rangle|\le N}} W_R(F).
\]
The inner minimum grants the optimal derivation for free; search cost is not charged (this is the user's stated convention and it is sound for a lower bound on executing a derivation).
We call $R$ \emph{size-faithful} if there is a polynomial $q$ with $\operatorname{cost}_R(\pi)\le|\pi|\le q(\operatorname{cost}_R(\pi)+|F|)$ for every valid $\pi$. All unit-weight DAG calculi in the corpus with polynomial-size nodes are size-faithful.
\begin{align*}
\mathrm{LB}(R)&:\iff \forall k,C\ \exists^\infty N:\ \overline W_R(N)>C(N+1)^k,\\
\mathrm{B3}(R)&:\iff \bigl(\mathrm{SAT}\in\mathrm{P}\implies \exists C,k\ \forall N:\ \overline W_R(N)\le C(N+1)^k\bigr).
\end{align*}
$\mathrm{LB}(R)\wedge\mathrm{B3}(R)\Rightarrow \mathrm{P}\ne\mathrm{NP}$ is immediate. The question of this note: for \emph{which} $R$ could B3 be a genuinely weaker, independently provable statement?

\section{Exact structural facts}
\begin{proposition}[Opposite monotonicity]\label{closure}
Let $R,R'$ be size-faithful and let $R'$ p-simulate $R$ in proof length. Then $\mathrm{B3}(R)\Rightarrow\mathrm{B3}(R')$ and $\mathrm{LB}(R')\Rightarrow\mathrm{LB}(R)$.
\end{proposition}
\begin{proof}
Simulation gives $\overline W_{R'}(N)\le r(\overline W_{R}(N)+N)$ for a polynomial $r$ (use faithfulness on both sides). A polynomial bound on $\overline W_R$ therefore transfers up; a superpolynomial bound on $\overline W_{R'}$ transfers down.
\end{proof}
So LB is closed \emph{downward} and B3 \emph{upward} in the simulation order. A proof of P$\ne$NP along this route needs a single system in the intersection.

\begin{proposition}[Conservation of difficulty]\label{conservation}
If $\mathrm{LB}(R)$ is a proved theorem, then $\mathrm{B3}(R)\iff \mathrm{P}\neq\mathrm{NP}$.
\end{proposition}
\begin{proof}
($\Rightarrow$) $\mathrm{B3}(R)$ and $\mathrm{LB}(R)$ give $\mathrm{SAT}\notin\mathrm{P}$. ($\Leftarrow$) If $\mathrm{SAT}\notin\mathrm{P}$, the implication defining B3 is vacuously true.
\end{proof}
\begin{corollary}\label{weak}
For resolution (Haken 1985), cutting planes (Pudl\'ak 1997), polynomial calculus (Razborov 1998), bounded-depth Frege (Ajtai 1988; Beame et al., Kraj\'{\i}\v{c}ek et al.\ 1993), and every system p-simulated by one of them (Proposition~\ref{closure}), proving B3 \emph{is} proving P$\ne$NP. For these systems there is no weaker bridge: B3 carries the whole difficulty.
\end{corollary}
This applies to the user's multisign DAG hypergraph reduction whenever its moves are local signed-hyperedge merges, unit propagation, or subsumption. Each such move is a constant number of resolution or weakening steps, so resolution p-simulates the system. Such a system's minimal reduced state can be shown exponentially costly on PHP or Tseitin families, but then the bridge is equivalent to the target itself.

\begin{proposition}[World classification]\label{world}
In Impagliazzo's Heuristica, Pessiland, Minicrypt and Cryptomania (all with $\mathrm{P}\ne\mathrm{NP}$), $\mathrm{B3}(R)$ is true for every $R$, vacuously. In Algorithmica it is equivalent to ``$R$ is polynomially bounded''. Hence every proof of $\mathrm{B3}(R)$ is, in effect, an argument inside Algorithmica that constructs short $R$-refutations from a hypothetical polynomial-time SAT decider.
\end{proposition}
\begin{proof}Immediate from the definition.\end{proof}
\begin{remark}
In the project's world ledger (\texttt{research/agents/claude/world-ledger/}) B3 would be classified as an \emph{implication} whose content lives only in Algorithmica, not a world separator. Proposed rows: B3(Res) $\equiv$ P$\ne$NP; B3(optimal $R$): implied by optimality (already in the corpus as LSPV-1/336); B3(EF): split below.
\end{remark}

\section{The only known contentful bridge: provable correctness}
Let $\mathrm{Corr}_A$ be the $\forall\Pi^b_1$ statement
$\forall x\,\forall y\,[\,A(x)=0\wedge |y|\le|x|\to \neg\mathrm{Sat}(x,y)\,]$, where $A$ is a clocked polynomial-time machine, formalized in Cook's theory PV.

\begin{theorem}[KNOWN THEOREM; Cook 1975, cf.\ Kraj\'{\i}\v{c}ek 1995 Ch.~9, Cook--Kraj\'{\i}\v{c}ek 2007]\label{pvbridge}
If $A$ is a polynomial-time SAT decider with $\mathrm{PV}\vdash\mathrm{Corr}_A$, then Extended Frege is polynomially bounded on CNF refutations: $\overline W_{\mathrm{EF}}(N)\le \mathrm{poly}(N)$.
\end{theorem}
\begin{proof}[Proof sketch]
Cook's translation gives polynomial-size EF proofs of the tautologies $\llbracket\mathrm{Corr}_A\rrbracket_n$ in free variables $\vec x,\vec y$, with extension variables for the gates of the circuit computing $A$ at length $n$. Substitute the bits of a concrete UNSAT $F$ for $\vec x$. EF verifies the constant-input evaluation $A(F)=0$ gate by gate in polynomial size. Modus ponens then yields a polynomial-size EF proof of $\neg\mathrm{Sat}(F,\vec y)$, i.e.\ of $\neg F(\vec y)$, which converts to an EF refutation of $F$. Cook--Kraj\'{\i}\v{c}ek (JSL 72, 2007) prove the sharper equivalence: $\mathrm{PV}\vdash \mathrm{P}=\mathrm{NP}$ iff PV defines a polynomial-time function that outputs, for every formula, either a satisfying assignment or an EF refutation.
\end{proof}

\begin{definition}[Residual gap]
$\mathrm{RG}(\mathrm{EF})$: if $\mathrm{SAT}\in\mathrm{P}$, then $\overline W_{\mathrm{EF}}$ is polynomially bounded \emph{even if} no polynomial-time SAT decider has a PV-provable $\mathrm{Corr}_A$.
\end{definition}

\begin{corollary}[Exact split of B3(EF)]\label{split}
$\mathrm{B3}(\mathrm{EF})$ follows from the conjunction of (i) Theorem~\ref{pvbridge} (KNOWN) and (ii) either $\mathrm{RG}(\mathrm{EF})$ or the reflection hypothesis
\[
\mathrm{RH}_{\mathrm{PV}}:\quad \mathrm{SAT}\in\mathrm{P}\implies \exists A\ \text{poly-time with}\ \mathrm{PV}\vdash \mathrm{Corr}_A .
\]
Consequently $\mathrm{LB}(\mathrm{EF})\wedge\mathrm{RH}_{\mathrm{PV}}\Rightarrow\mathrm{P}\ne\mathrm{NP}$. Both conjuncts are OPEN.
\end{corollary}
\begin{proof}
Under SAT$\in$P, $\mathrm{RH}_{\mathrm{PV}}$ supplies an $A$ satisfying the hypothesis of Theorem~\ref{pvbridge}. $\mathrm{RG}(\mathrm{EF})$ gives the conclusion directly in the remaining case.
\end{proof}
The same pattern lifts to any pair (theory $T$, propositional system $R_T$) with a Cook-style translation. Examples include $T^1_2/\mathrm{S}^2_2$ with bounded-depth $G$ systems, and $T$ with $R_T=\mathrm{EF}+\llbracket\mathrm{Ref}_T\rrbracket$. Strengthening $T$ weakens $\mathrm{RH}_T$ but moves $R_T$ \emph{up} the simulation order, where LB becomes harder (Proposition~\ref{closure}). No fixed sound r.e.\ $T$ proves $\mathrm{Corr}_A$ for every correct clocked $A$ (corpus: LSPV-1/336, ``No effective enumeration''). Choosing $T$ after $A$ is therefore not allowed for a preselected $R$. Optimality of $R$ (LSPV-1/336) remains the only known way to remove the dependence on $T$.

\section{Consequences for the five assignments}
\begin{enumerate}
\item \textbf{Lead / Builder.} If the multisign hypergraph calculus is to carry a contentful B3, it must reach at least the EF level, for example by letting a move introduce a new DAG node defined as a gate over existing nodes (extension). At that level Theorem~\ref{pvbridge} supplies the PV portion of the bridge for free. The open obligations become $\mathrm{RH}_{\mathrm{PV}}$ (or $\mathrm{RG}$) and $\mathrm{LB}(\mathrm{EF})$.
\item \textbf{Hard families.} PHP and Tseitin, which are hard for resolution, have polynomial-size EF (indeed Frege) refutations (Cook--Reckhow 1979; Buss 1987; Urquhart). This is a FAMILY FAILURE at the EF level. Random $\Delta n$-clause 3-CNF and Kraj\'{\i}\v{c}ek-style proof-complexity generators remain CANDIDATE EF-hard families; no superpolynomial EF lower bound is known for any family.
\item \textbf{Resource auditor.} Theorem~\ref{pvbridge} is pointwise in the sense of B1: $W_{\mathrm{EF}}(F)\le p(|F|+T_A(F))$ when $\mathrm{Corr}_A$ is PV-provable, with $p$ depending on the size of the PV proof. This is a correct-quantifier instance of B1 for PV-verified algorithms. It does not contradict the corpus counterexample to B1 for resolution.
\item \textbf{Red team.} Pich--Santhanam (STOC 2021) prove that PV cannot prove strong co-nondeterministic circuit lower bounds for NP. ``Why are proof complexity lower bounds hard?'' (FOCS 2019) shows that at least one of two obstacles holds for EF. Together these are barrier evidence that $\mathrm{LB}(\mathrm{EF})$ is hard to obtain by feasible reasoning. Any claimed hypergraph EF lower bound should be checked against them. Classification: directly transferable barrier theorems, not analogies.
\end{enumerate}

\section*{CORPUS REVIEWED}
\texttt{AGENTS.md}; the 2026-11-02 worst-case directive and LSPV-1 core-barrier handoff 335; LSPV-1/336 (size-optimality bridge, non-r.e.\ correctness index); the 2026-11-13 TMR replacement directive; the 2026-11-15 Allender integration directive; the recent TMR-1..5 commit log; the \texttt{research/agents/claude/} world ledger. A repository search for Cook--Kraj\'{\i}\v{c}ek / PV-to-EF found only bibliography mentions, with no prior use as a B3 split.

\section*{DUPLICATION / CONTRADICTION AUDIT}
The optimality sufficient bridge and the non-r.e.\ correctness barrier are \emph{adopted} from LSPV-1/336, not new. New in this note: opposite monotonicity, the conservation equivalence for proved-LB systems, the world classification of B3, and the PV/RG split of B3(EF). No contradiction with root \texttt{main.tex} was found. This note does not touch the TMR Allender priority.

\section*{WHAT CHANGED RELATIVE TO ROOT MAIN}
Nothing in root \texttt{main.tex} is edited. This routine is not an LSPV/TMR role and owns no region there. The note is offered as a supporting artifact.

\section*{ADOPTED / REJECTED / RETIRED PRIOR COMPONENTS}
Adopted: LSPV-1/336. Rejected as a contentful route: B3 for any system p-simulated by resolution, CP, PC or bounded-depth Frege (Corollary~\ref{weak}). These remain valid as LB testbeds only.

\section*{STALE OR OVERSTATED CLAIMS FOUND}
None detected. Caution for future runs: ``prove the weakest bridge B3 for the multisign hypergraph calculus'' is \emph{not} weaker than P$\ne$NP when that calculus is resolution-simulable.

\section*{References}
\begin{itemize}
\item S.~Cook, Feasibly constructive proofs and the propositional calculus, STOC 1975.
\item S.~Cook, R.~Reckhow, The relative efficiency of propositional proof systems, JSL 44 (1979).
\item S.~Cook, J.~Kraj\'{\i}\v{c}ek, Consequences of the provability of NP $\subseteq$ P/poly, JSL 72 (2007).
\item J.~Kraj\'{\i}\v{c}ek, Bounded Arithmetic, Propositional Logic and Complexity Theory, CUP 1995.
\item J.~Kraj\'{\i}\v{c}ek, P.~Pudl\'ak, Propositional proof systems, the consistency of first order theories and the complexity of computations, JSL 54 (1989).
\item J.~Pich, R.~Santhanam, Strong co-nondeterministic lower bounds for NP cannot be proved feasibly, STOC 2021. \url{https://ora.ox.ac.uk/objects/uuid:4bb0ed7a-1983-4d95-8296-ab973094292e}
\item J.~Pich, R.~Santhanam, Why are proof complexity lower bounds hard?, FOCS 2019.
\item Unprovability of strong complexity lower bounds in bounded arithmetic, \url{https://arxiv.org/pdf/2305.15235}.
\end{itemize}

\noindent P versus NP remains OPEN.
\end{document}
% END_OF_VERIFIED_REPORT::CLAUDE::B3-FRONTIER-0001
